\scanpage{363}\mspage{27\deleted{\unsure{8}}9}
Posons \deleted{\unsure{donc}}
\[
 \pi'=2^\pi,
\]
on a donc ${\pi'}^{\pi_0}=2^{\pi\pi_0}=2^\pi=\pi'$, et \deleted{\unsure{on a}}
$\operatorname{Card\,red}M_\pi\leq\pi'$. Donc 4.9.5. nous donne
que
\[
 \underline{\mathrm{Hom}}(M_\pi,N)\simeq\mathrm{Ind}_{\pi'}(\mathrm{Hom}(M_{\unsure{\pi}},M_{\pi'})),
\]
d'où

\uline{Corollaire} 4.12.5. Sous les conditions de
4.12.4, si $\pi_0$ est utile pour $M$ \uline{et}
pour $N$, alors pour $\pi\geq\pi_1$, \deleted{\unsure{on aura}}
$\underline{\mathrm{Hom}}_{!\pi}(M,N)$ est \deleted{$\unsure{\mathrm I}$}$\pi'$-accessible, où $\pi'=2^\pi$,
de façon plus précise on a une
\unsure{équivalence} de catégories
\[
 \underline{\mathrm{Hom}}_{!\pi}(M,N)\unsure{\underline{\approx}|\xrightarrow{\approx}}
 \mathrm{Ind}_{\pi'}(\underline{\mathrm{Hom}}(M_\pi,N_{\pi'})).
\]

\uline{Remarque} 4.12.6. Si \deleted{\unsure{cat}} $M,N$ sont accessible\unsure{s},
la catégorie $\underline{\mathrm{Hom}}_{\mathrm{acc}}(M,N)$ des foncteurs
accessibles n'est pas \deleted{\unsure{forcément}} tjrs accessible.
Si $N$ satisfait $L_{\pi_0}$, cette catégorie
satisfait $L_{\pi_0}$, et les $\unsure{\underrightarrow{\lim}|\underline{\lim}}$ gdes devant
$\pi_0$ se calculent argument par argument.
% archive sheet number: 362

